% Tier Two results sections (A1-A4), mirroring the main-text format with T2 numbers.
% Sources: publication/forest_plots/total_adjusted/t2_sorted/*_data.csv
% Corrected significance approximated by z = |log RR| / se > 3.19 (alpha = .05/34).

\subsection{Effect of the overall availability of alternative proteins on general meat purchases (A1, Tier Two)}

For analysis set A1, we report all estimates as the percentage change in purchases of the outcome, relative to total purchases, associated with a 10-percentage-point (proportion) or one-menu-item (count) increase in availability. We found little evidence that any of the overall availability exposures affected general ABF purchases: with one exception, pooled point estimates were all within 15 percentage points of null. The exception, vegetarian menu availability (as a proportion), was not significant upon multiplicity correction; its point estimate suggested a reduction in chicken- and fish-containing purchases by -17\% (95\% CI: -26\%, -6.1\%). The same exposure was associated with a reduction in meat-containing purchases by -9.9\% (95\% CI: -17\%, -1.8\%); these were the only two associations significant without correction, and neither was significant after correction.

For most exposure-outcome pairings in A1, within-restaurant estimates varied in direction across restaurants, and no pairing had every restaurant's estimate in the expected direction. The most consistent pairings were vegetarian availability (as a proportion) on meat purchases (16 of 18 restaurants in the expected direction) and on chicken- and fish-containing purchases (15 of 17). The largest estimates in the hypothesized direction came from Restaurants 6 (coffee shop) and 7 (juice bar): a 10-percentage-point increase in vegetarian-menu share was associated with respective relative shifts of -37\% in chicken and fish purchases (95\% CI: -49\%, -23\%) and -32\% in meat purchases (95\% CI: -42\%, -20\%).

\subsection{Effect of the overall availability of alternative proteins on counterpart-specific meat purchases (A2, Tier Two)}

For analysis set A2, we report all estimates as the percentage change in purchases of the counterpart-specific outcomes, relative to total purchases, associated with each additional alternative protein emulating that counterpart on the menu (``count''); presence exposures are omitted for Tier Two. Across all five protein product classes used in this analysis set, we found no evidence that the alternative protein exposures reduced purchases of their counterpart ABF products; all five pooled point estimates were in the unexpected direction. One association, ground meat, was significant without correction (+20\%; 95\% CI: +0.5\%, +82\%), and none were significant after correction.

Within-restaurant estimates for count exposures reached uncorrected significance in about half of cases (21 of 43), nearly all in the unexpected direction, and ten survived correction. The only corrected-significant estimate in the expected direction was breakfast-style meat at Restaurant 20 (halfsmoke and chili shop), at -19\%. The largest estimates came from low-volume restaurants and should be interpreted with the purchase counts in mind.

\subsection{Effect of the introduction of new alternative proteins on general meat purchases (A3, Tier Two)}

For analysis set A3, we report all estimates as the percentage change in purchases of the outcome, relative to total purchases, associated with the introduction of a new alternative protein, as an immediate change at launch (``level change'') and a change in the annualized trend (``slope change''). New introductions of alternative proteins produced no changes in general ABF purchases at launch or over the following year at the pooled level, even at the uncorrected 95\% level. All pooled level change estimates were within 19 percentage points of null, and four of five were in the expected direction. Slope-change estimates were more dispersed due to higher statistical uncertainty; three of five were in the unexpected direction, though this should be interpreted cautiously given the wide CIs.

Within-restaurant level-change estimates reached uncorrected significance in 27 of 109 cases, of which 12 survived correction, concentrated in a small number of restaurants. Two-thirds of estimates were in the expected direction, and most were within 50 percentage points of the null; the largest deviations came from one low-volume restaurant with small purchase counts. Within-restaurant slope-change estimates were too noisy to usefully interpret.

\subsection{Effect of the introduction of alternative proteins on counterpart-specific purchases (A4, Tier Two)}

For analysis set A4, we report all pooled point estimates as the percentage change in counterpart-specific meat purchases, relative to total purchases, at launch and over the subsequent year associated with the introduction of an alternative protein emulating that counterpart. Newly introduced counterpart-specific alternative proteins produced no significant changes in counterpart purchases, even at the uncorrected 95\% level. Pooled level change estimates were mixed: breakfast-style meat declined while chicken, dairy, and ground meat rose, none significantly. Pooled slope change estimates were directionally consistent, all four in the expected direction.

Within-restaurant estimates reached uncorrected significance in 17 of 32 cases, of which six survived correction, four of these in the unexpected direction; among them was the whole-muscle meat estimate at Restaurant 1 (Greek rotisserie chain), an outcome shown at the restaurant level only. Three-quarters of level-change estimates were within roughly 55 percentage points of the null; the remainder were implausibly large values from restaurants with sparse purchase counts. As before, the slope-change estimates should be interpreted cautiously due to statistical noise.
